\paragraph{}To demonstrate how power factor correction works we will
set up two different circuits.
The first will contain a purely resistive load of 600$\Omega$ and
should therefore have a power factor of 1.
And the second circuit will have a complex load with an inductive
impedance of 300+j300$\Omega$.
With each circuit we will connect a parallel capacitive impedance of
8 specific values, -[171, 200, 240, 300, 400, 600, 1200]$\Omega$.
During the build of the circuit we had an easier time getting
everything connected than we did  the previous lab. One main point of
improvement came when switching between the two circuits. For this we
learned how we could essentially shift all connection up by one,
making the load changes easier.

\paragraph{}With both circuits we measured the line voltage, phase
voltage, line current, the real and reactive power,
and the power factor. Samples were captured for all 8 capacitive
loads. In addition, we calculated theoretical
values for all the same loads. We also calculated these values by
hand using the equations listed bellow in section 2.1.
This allowed us to compare the measured values to expected values.

\clearpage
\subsection{Methods of Calculation}
\begin{itemize}
  \item $\bar{Z}_\|$ - Resulting parallel impedance after adding the
    capacitive load
    \begin{center}
      $\bar{Z}_{\|} = \frac{
        \bar{Z}_{1}\bar{Z}_{2}
      }{
        \bar{Z}_1 + \bar{Z}_2
      }$
    \end{center}

  \item $\bar{I}_L$ - Line current in the phase
    \begin{center}
      $\bar{I}_{L} = \frac{
        V_{\phi}
      }{
        \bar{Z_{\|}}
      }$
    \end{center}

  \item $\bar{S}$ - Complex power in VA
    \begin{center}
      $\bar{S} = 3V_{\phi}I_{L}^*$
    \end{center}

  \item $P$ - Real power in W
    \begin{center}
      $P = \Re(\bar{S})$
    \end{center}

  \item $Q$ - Reactive power in var
    \begin{center}
      $Q = \Im(\bar{S})$
    \end{center}

  \item $\theta$ - Power factor angle
    \begin{center}
      $\theta = \arg(\bar{S})$
    \end{center}

  \item PF - Power factor
    \begin{center}
      PF = $\cos(\theta)$
    \end{center}

  \item PF$_{c}$ - Power factor calculated from measured P, Q
    \begin{center}
      PF$_{c}  = \frac{
        P_{m}
      }{
        |P_{m} + jQ_{m}|
      }$
    \end{center}

  \item $\theta_{m}$ - Power factor angle calculated from the
    measured power factor
    \begin{center}
      $\theta_{m} = \arccos($PF$_{m})$
    \end{center}

  \item
    $\theta_{c}$ - Power factor angle derived from measured $P$ and $Q$
    \begin{center}
      $\theta_{c} = \arctan2(P, Q)$
    \end{center}

\end{itemize}

\clearpage
